%!TEX TS-program = xelatex
%!TEX encoding = UTF-8 Unicode
\documentclass[10pt,a4paper,withhyper,ragged2e]{altacv}

\geometry{left=1.15cm,right=1.15cm,top=1.15cm,bottom=1.2cm,columnsep=0.95cm}
\usepackage{paracol}
\usepackage{xeCJK}
\usepackage{qrcode}

\setmainfont{Source Sans 3}
\setsansfont{Source Sans 3}
\setCJKmainfont{Microsoft YaHei}
\setCJKsansfont{Microsoft YaHei}
\renewcommand{\familydefault}{\sfdefault}

\definecolor{AcademicNavy}{HTML}{17324D}
\definecolor{AcademicTeal}{HTML}{0B7A75}
\definecolor{AcademicGold}{HTML}{C7A54B}
\definecolor{BodyGrey}{HTML}{4E5964}
\definecolor{SoftGrey}{HTML}{74808B}
\colorlet{name}{AcademicNavy}
\colorlet{tagline}{AcademicTeal}
\colorlet{heading}{AcademicNavy}
\colorlet{headingrule}{AcademicGold}
\colorlet{subheading}{AcademicTeal}
\colorlet{accent}{AcademicTeal}
\colorlet{emphasis}{AcademicNavy}
\colorlet{body}{BodyGrey}

\renewcommand{\namefont}{\Huge\bfseries}
\renewcommand{\taglinefont}{\large\bfseries}
\renewcommand{\personalinfofont}{\small}
\renewcommand{\cvsectionfont}{\Large\bfseries}
\renewcommand{\cvsubsectionfont}{\large\bfseries}
\renewcommand{\cvItemMarker}{{\small\faAngleRight}}
\renewcommand{\cvRatingMarker}{\faCircle}
\renewcommand{\locationname}{地点}
\renewcommand{\datename}{时间}
\setlength{\parindent}{0pt}

\begin{document}

\name{马睿骁\quad {\LARGE Ruixiao Ma}}
\tagline{物理学硕士研究生\enskip\textperiodcentered\enskip 材料计算、原子自旋模拟与机器学习}
\personalinfo{%
  \email{2024210104023@mails.zstu.edu.cn}
  \phone{微信：19800353973}
  \homepage{marx2001.github.io}
  \github{marx2001}
  \location{浙江\enskip 杭州}
}
\makecvheader

\columnratio{0.64}
\begin{paracol}{2}

\cvsection{研究概述\enskip Profile}

物理学硕士研究生，研究兴趣集中于二维磁性与多铁材料、拓扑物态及数据驱动的参数空间探索。主要使用第一性原理计算、Wannier 紧束缚模型和 Python 科研脚本开展材料物性分析，并持续整理可复现的计算流程与研究记录。

\cvsection{研究方向\enskip Research Interests}

\cvevent{二维多铁材料的第一性原理计算}{磁性\enskip\textperiodcentered\enskip 电性\enskip\textperiodcentered\enskip 谷电子学}{2024.09 -- 至今}{浙江理工大学}
\begin{itemize}
  \item \textbf{计算材料：}研究二维磁性、多铁性及谷相关物理性质。
  \item 结合能带、态密度与磁交换参数，分析结构和电子性质之间的关联。
\end{itemize}

\divider

\cvevent{机器学习}{拓扑相界\enskip\textperiodcentered\enskip 拓扑绝缘体}{2025 -- 至今}{浙江理工大学}
\begin{itemize}
  \item \textbf{拓扑模型：}结合 Wannier 化与紧束缚模型分析能带演化和拓扑表征。
  \item \textbf{数据驱动探索：}使用机器学习与参数空间扫描寻找磁性体系中的候选拓扑物态及相界。
\end{itemize}

\cvsection{研究与工作经历\enskip Experience}

\cvevent{计算物理课题组负责人}{浙江理工大学}{2025.09 -- 2026.09}{杭州}
\begin{itemize}
  \item 负责计算物理课题组开题答辩等会议材料的收集、整理与主持。
  \item 定期汇总计算物理课题组成员的科研与生活状态。
\end{itemize}

\divider

\cvevent{科研助理}{中国信息通信研究院\enskip 云计算与大数据研究所}{2022.02 -- 2022.05}{北京}
\begin{itemize}
  \item 参与机器学习方向的科研辅助工作，获得实习证明。
  \item 承担演示文稿制作、论文写作辅助及文字校对等工作。
\end{itemize}

\cvsection{语言能力\enskip Language}

\cvtag{大学英语六级 2021}\cvtag{大学英语四级 2020}
\cvtag{三一口语五级}

\cvsection{国际交流\enskip International Exchange}

\cvevent{俄罗斯、美国与英国游学交流}{院校访问与文化交流}{2015 -- 2017}{国际交流}
\begin{itemize}
  \item 参加俄罗斯游学与乐团交流演出，在美国访学交流，访问哈佛大学、耶鲁大学等院校。
  \item 在英国参加游学交流，访问牛津大学、剑桥大学、伦敦帝国理工学院等院校。
\end{itemize}

\switchcolumn

\cvsection{教育背景\enskip Education}

\cvevent{理学硕士\enskip 物理学}{浙江理工大学}{2024.09 -- 至今}{杭州}
\small 浙江省量子态调控与光场操控重点实验室\par
量子物态调控研究所

\divider

\cvevent{理学学士}{北京工商大学}{2020.09 -- 2024.07}{北京}
\small 光电信息科学与工程

\cvsection{学术论文\enskip Publications}

{\footnotesize
\textbf{Ferroelectric Control of Magnetism, Valley Polarization, and Skyrmions in Monolayer TcIrGe$_2$Se$_6$}\par
\textbf{Ruixiao Ma}, Zhangtong Ying, Zirui Zhang, et al.\par
\textcolor{accent}{\textbf{第一作者\enskip\textperiodcentered\enskip Physical Review B\enskip\textperiodcentered\enskip 已接收}}\par
}

\divider

{\footnotesize
\textbf{Pt-Doping-Induced Skyrmion Lattice in a Two-Dimensional Multiferroic Monolayer Beyond Room Temperature}\par
\textbf{Ruixiao Ma}, Zheng Chen, Zhen Sun, et al.\par
\textcolor{accent}{\textbf{第一作者\enskip\textperiodcentered\enskip JPCC\enskip\textperiodcentered\enskip Under review}}\par
}

\divider

{\footnotesize
\textbf{Interpretable Machine Learning of Spin-Chern Phases in Two-Dimensional Altermagnetic Tessellations}\par
\textbf{Ruixiao Ma}, Zhen Sun, Xiaoping Wu, Changsheng Song\par
\textcolor{accent}{\textbf{第一作者\enskip\textperiodcentered\enskip 修改中}}\par
}

\divider

{\footnotesize
\textbf{Room-Temperature Triferroic-Topological Coupling and All-Electric Meron-Skyrmion Switching in an AlCrSe$_3$/RuBr$_2$ van der Waals Heterostructure}\par
\textbf{Ruixiao Ma}, et al.\par
\textcolor{accent}{\textbf{第一作者\enskip\textperiodcentered\enskip 在写}}\par
}

\cvsection{荣誉与奖项\enskip Honors}

\cvachievement{\faAward}{学业三等奖学金}{浙江理工大学，8000 元，2026}

\divider

\cvachievement{\faAward}{学业二等奖学金}{浙江理工大学，10000 元，2025}

\divider

\cvachievement{\faGraduationCap}{新生学业奖学金}{浙江理工大学，8000 元，2024}

\divider

\cvachievement{\faTrophy}{北京市大学生艺术节金奖}{2023}

\divider

\cvachievement{\faMedal}{北京市三好学生}{2017}

\cvsection{兴趣与特长\enskip Interests}

\cvtag{钢琴}\cvtag{贝斯}\cvtag{棒球}
\cvtag{长跑}\cvtag{阅读}\cvtag{科研写作}

\end{paracol}

\newpage
\begin{paracol}{2}

\cvsection{实践经历\enskip Practice}

\cvevent{运营者与内容作者}{科普微信公众号“漫步科学世界”}{2020.09 -- 2024.07}{线上}
\begin{itemize}
  \item 围绕超材料课程学习成果、硅光芯片等主题开展科普内容创作。
  \item 硅光芯片相关文章阅读量超过一万；个人网站建成后暂停公众号运营。
\end{itemize}

\divider

\cvevent{寒暑假实习}{学而思}{2020.09 -- 2024.07}{北京}

\cvsection{校园活动与志愿服务\enskip Campus \& Volunteering}

\cvevent{贝斯首席}{北京工商大学校管乐团}{2020.09 -- 2024.07}{北京}
参与乐团活动，并获北京市大学生艺术节金奖。

\divider

\cvevent{干事与志愿者}{北京工商大学校青年志愿者协会}{2020.09 -- 2021.07}{北京}
参与建党 100 周年大会、北京房山国际长走大会及建校 70 周年志愿服务。

\divider

\cvevent{干事与赛事志愿者}{浙江理工大学}{2024.09 -- 2025.07}{杭州}
担任校联合会干事，并参与 2026 世界杯暨 2027 亚洲杯联合预选赛志愿服务。

\divider

\cvevent{贝斯首席}{德茂中学金帆民乐团}{2014.09 -- 2017.07}{北京}


\switchcolumn

\cvsection{学术记录\enskip Research Notes}

持续维护个人学术主页，系统整理第一性原理计算、软件学习、绘图流程、研究心得和代码记录。

\medskip
\begin{center}
  \qrcode[height=1.7cm]{https://marx2001.github.io/}\\[0.4em]
  \href{https://marx2001.github.io/}{\small marx2001.github.io}
\end{center}

\cvsection{研究工具\enskip Toolkit}

\cvtag{VASP}\cvtag{Wannier90}\cvtag{PythTB}\par\smallskip
\cvtag{Pybinding}\cvtag{IRVSP}\cvtag{Pyw90}\par\smallskip
\cvtag{Python}\cvtag{机器学习}\cvtag{LaTeX}\par\smallskip
\cvtag{Blender}

\cvsection{方法能力\enskip Methods}

\cvtag{第一性原理计算}\cvtag{紧束缚模型}
\cvtag{能带分析}\cvtag{态密度分析}
\cvtag{磁交换参数}\cvtag{拓扑性质分析}
\cvtag{科研流程整理}


\cvsection{个人信息\enskip Personal}

\cvachievement{\faCalendar}{出生年月}{2001.11}

\divider

\cvachievement{\faMapMarker}{籍贯}{北京大兴}

\divider

\cvachievement{\faFlag}{政治面貌}{共青团员}




\end{paracol}

\end{document}
